\documentclass[12pt]{article}
\usepackage[utf8]{inputenc}
\usepackage[T1]{fontenc}
\usepackage[margin=1in]{geometry}
\usepackage{setspace}
\usepackage{xcolor}
\usepackage[colorlinks=true, urlcolor=blue!60!black, linkcolor=blue!60!black]{hyperref}
\setlength{\parindent}{0em}
\setlength{\parskip}{0.45em}
\renewcommand{\baselinestretch}{1.0}
\begin{document}
\begin{center}
{\Large \textit{Puerto Rico v. Branstad}, 483 U.S. 219 (1987)\par}
\vspace{0.3em}
{\normalsize Case Brief --- One-Page Edition\par}
\end{center}

\textbf{Memory Jogger.} Overruling \textit{Kentucky v.\ Dennison}: federal courts can compel governors to extradite fugitives (\href{https://supreme.justia.com/cases/federal/us/483/219/}{483 U.S. 219}).

\textbf{Facts.} Ronald Calder, charged with murder in Puerto Rico for running over a pregnant woman, fled to Iowa; Iowa's Governors refused extradition, questioning whether he could get a fair trial.

\textbf{Procedural History.} The lower federal courts dismissed Puerto Rico's mandamus suit under \textit{Dennison}.

\textbf{Votes.} Reversed, 9--0 (6--3 on the constitutional holding): Marshall for the Court; O'Connor and Scalia concurred in part.

\textbf{Holding.} Federal courts may compel an asylum governor to perform the ministerial duty of extradition, and the Extradition Act gives Puerto Rico that right as a Territory.

\textbf{Reasoning.} \textit{Dennison}, decided on the eve of the Civil War, rested on a coequal-sovereigns view long since rejected; federal courts routinely enforce federal duties against state officers, and a duty imposed by the Constitution itself needs no Tenth Amendment balancing. O'Connor would rest on the statute alone; Scalia reserved Congress's power over Territory extradition.

\textbf{Examples.} \emph{Same-Side:} a shield law cannot protect an actual fugitive. \emph{Opposite-Side:} a telehealth prescriber who never fled. \emph{Fence-Sitter:} Puerto Rico's statutory right after \textit{Printz}.

\textbf{Critique.} Resting the Territory remedy on statute invites \textit{Printz} problems, and mandatory rendition removes governors' humanitarian discretion.

\textbf{Key Quotation.} ``\emph{Kentucky v. Dennison} is the product of another time. The conception of the relation between the States and the Federal Government there announced is fundamentally incompatible with more than a century of constitutional development. Yet this decision has stood while the world of which it was a part has passed away. We conclude that it may stand no longer.'' --- Justice Marshall, Opinion of the Court, 483 U.S. at 230.

\end{document}
